\textbf{Phần 1. Trắc nghiệm lựa chọn}
\setcounterpageref{ex}{0}

\begin{ex} %[1D6V24]
Tìm số tròn chục lớn nhất có hai chữ số.
\choice
{$80$}
{\True $90$}
{$100$}
{$89$}
\loigiai{
Số tròn chục lớn nhất có hai chữ số là số $90$.
\textbf{Đáp số:} $90$.\par
\textbf{Đáp án của câu hỏi là B}
}
\end{ex}

\begin{ex} %[1D6V24]
Tìm số lớn nhất có hai chữ số.
\choice
{\True $99$}
{$98$}
{$100$}
{$97$}
\loigiai{
Số lớn nhất có hai chữ số là số $99$.
\textbf{Đáp số:} $99$.\par
\textbf{Đáp án của câu hỏi là A}
}
\end{ex}

\begin{ex} %[1D6V24]
Tìm số liền sau của số lớn nhất có hai chữ số.
\choice
{$90$}
{$110$}
{\True $100$}
{$99$}
\loigiai{
Số liền sau của 99 là số $100$.
\textbf{Đáp số:} $100$.\par
\textbf{Đáp án của câu hỏi là C}
}
\end{ex}

\begin{ex} %[1D6V24]
Số 75 gồm mấy chục và mấy đơn vị? Lấy chữ số hàng chục trừ đi chữ số hàng đơn vị được kết quả là bao nhiêu?
\choice
{$1$}
{$3$}
{$4$}
{\True $2$}
\loigiai{
Số 75 có chữ số hàng chục là 7, hàng đơn vị là 5: $7 - 5 = 2$.
\textbf{Đáp số:} $2$.\par
\textbf{Đáp án của câu hỏi là D}
}
\end{ex}

\begin{ex} %[1D6V24]
Cho ba chữ số: 3, 0, 8. Lập số lớn nhất có hai chữ số từ các chữ số đã cho.
\choice
{$82$}
{\True $83$}
{$84$}
{$81$}
\loigiai{
Chọn hàng chục lớn nhất là 8, hàng đơn vị lớn nhất là 3: số đó là $83$.
\textbf{Đáp số:} $83$.\par
\textbf{Đáp án của câu hỏi là B}
}
\end{ex}

\begin{ex} %[1D6V24]
Cho ba chữ số: 3, 0, 8. Lập số bé nhất có hai chữ số khác nhau từ các chữ số đã cho.
\choice
{\True $30$}
{$20$}
{$40$}
{$29$}
\loigiai{
Hàng chục bé nhất khác 0 là 3, hàng đơn vị bé nhất là 0: số đó là $30$.
\textbf{Đáp số:} $30$.\par
\textbf{Đáp án của câu hỏi là A}
}
\end{ex}

\begin{ex} %[1D6V24]
Cho các số: $56, 65, 59, 62$. Sắp xếp theo thứ tự từ bé đến lớn, số đứng ở vị trí thứ ba là số nào?
\choice
{$61$}
{$63$}
{\True $62$}
{$60$}
\loigiai{
Thứ tự từ bé đến lớn: $56 < 59 < 62 < 65$. Số đứng thứ ba là số $62$.
\textbf{Đáp số:} $62$.\par
\textbf{Đáp án của câu hỏi là C}
}
\end{ex}

\begin{ex} %[1D6V24]
Tìm số tự nhiên $a$ biết rằng: $69 < a < 71$.
\choice
{$60$}
{$80$}
{$69$}
{\True $70$}
\loigiai{
Số tự nhiên nằm giữa 69 và 71 là số tròn chục $70$.
\textbf{Đáp số:} $70$.\par
\textbf{Đáp án của câu hỏi là D}
}
\end{ex}

\begin{ex} %[1D6V24]
Bác nông dân có 5 chục quả trứng gà và 3 chục quả trứng vịt. Hỏi bác nông dân có tất cả bao nhiêu quả trứng cả hai loại?
\choice
{$70$}
{\True $80$}
{$90$}
{$79$}
\loigiai{
Tổng số chục trứng là: $5 + 3 = 8$ (chục).\n8 chục quả trứng tương ứng với $80$ quả trứng.
\textbf{Đáp số:} $80$.\par
\textbf{Đáp án của câu hỏi là B}
}
\end{ex}

\begin{ex} %[1D6V24]
Một cửa hàng có 90 chiếc cặp sách, đã bán được 40 chiếc cặp sách. Hỏi cửa hàng còn lại bao nhiêu chiếc cặp sách?
\choice
{\True $50$}
{$40$}
{$60$}
{$49$}
\loigiai{
Số cặp sách còn lại là: $90 - 40 = 50$ (chiếc cặp sách).
\textbf{Đáp số:} $50$.\par
\textbf{Đáp án của câu hỏi là A}
}
\end{ex}

\begin{ex} %[1D6V24]
Tìm số tròn chục $x$ thỏa mãn: $20 + x = 70$. Giá trị của $x$ là bao nhiêu?
\choice
{$40$}
{$60$}
{\True $50$}
{$49$}
\loigiai{
Ta có: $x = 70 - 20 = 50$.
\textbf{Đáp số:} $50$.\par
\textbf{Đáp án của câu hỏi là C}
}
\end{ex}

\begin{ex} %[1D6V24]
Tìm số tròn chục $y$ thỏa mãn: $80 - y = 30$. Giá trị của $y$ là bao nhiêu?
\choice
{$40$}
{$60$}
{$49$}
{\True $50$}
\loigiai{
Ta có: $y = 80 - 30 = 50$.
\textbf{Đáp số:} $50$.\par
\textbf{Đáp án của câu hỏi là D}
}
\end{ex}

\textbf{Phần 2. Trắc nghiệm đúng sai}
\setcounterpageref{ex}{0}

\begin{ex} %[1D6V24]
Xét tính đúng hoặc sai của các khẳng định toán học sau đây:
\choiceTF
{Bạn Lan có 35 que tính, bạn Mai có 38 que tính. bạn Mai có nhiều hơn bạn Lan $2$ que tính.}
{\True Lớp 1A có 32 học sinh, lớp 1B có 35 học sinh. cả hai lớp có tất cả $67$ học sinh.}
{\True Kết quả của bài toán «Tìm số có hai chữ số giống nhau biết tổng hai chữ số bằng 8» là $44$.}
{Kết quả của bài toán «Tìm số có hai chữ số giống nhau biết tổng hai chữ số bằng 12» là $67$.}
\loigiai{
\begin{itemize}
    \item Ý a) Sai: Kết quả đúng là $3$, không phải $2$. Số que tính Mai nhiều hơn Lan là: $38 - 35 = 3$ (que tính).
\textbf{Đáp số:} $3$.
    \item Ý b) Đúng: Tổng số học sinh cả hai lớp là: $32 + 35 = 67$ (học sinh).
\textbf{Đáp số:} $67$.
    \item Ý c) Đúng: Hai chữ số giống nhau có tổng bằng 8 là 4 và 4. Số đó là $44$.
\textbf{Đáp số:} $44$.
    \item Ý d) Sai: Kết quả đúng là $66$, không phải $67$. Hai chữ số giống nhau có tổng bằng 12 là 6 và 6. Số đó là $66$.
\textbf{Đáp số:} $66$.
\end{itemize}
\begin{itemize}
    \item a) \textbf{Sai}.
    \item b) \textbf{Đúng}.
    \item c) \textbf{Đúng}.
    \item d) \textbf{Sai}.
\end{itemize}
}
\end{ex}

\begin{ex} %[1D6V24]
Xét tính đúng hoặc sai của các bài toán và phép tính sau đây:
\choiceTF
{\True Kết quả của bài toán «Hiệu của số tròn chục lớn nhất và số tròn chục nhỏ nhất có hai chữ số là bao nhiêu» là $80$.}
{Kết quả của bài toán «Từ số 40 đến số 50 có tất cả bao nhiêu số tự nhiên liên tiếp» là $10$.}
{\True Một người có 60 bông hoa, người đó cắm vào các lọ mỗi lọ 10 bông hoa. cắm được tất cả $6$ lọ hoa.}
{Kết quả của bài toán «Tìm số có hai chữ số giống nhau biết tổng hai chữ số bằng 8» là $45$.}
\loigiai{
\begin{itemize}
    \item Ý a) Đúng: Số tròn chục lớn nhất là 90, nhỏ nhất là 10: $90 - 10 = 80$.
\textbf{Đáp số:} $80$.
    \item Ý b) Sai: Kết quả đúng là $11$, không phải $10$. Số lượng số là: $50 - 40 + 1 = 11$ (số).
\textbf{Đáp số:} $11$.
    \item Ý c) Đúng: 60 bông hoa gồm 6 chục bông hoa, chia mỗi lọ 1 chục nên cắm được $6$ lọ hoa.
\textbf{Đáp số:} $6$.
    \item Ý d) Sai: Kết quả đúng là $44$, không phải $45$. Hai chữ số giống nhau có tổng bằng 8 là 4 và 4. Số đó là $44$.
\textbf{Đáp số:} $44$.
\end{itemize}
\begin{itemize}
    \item a) \textbf{Đúng}.
    \item b) \textbf{Sai}.
    \item c) \textbf{Đúng}.
    \item d) \textbf{Sai}.
\end{itemize}
}
\end{ex}

\textbf{Phần 3. Trả lời ngắn}
\setcounterpageref{ex}{0}

\begin{ex} %[1D6V24]
Đoàn tàu hỏa có 10 toa xe. Toa thứ 7 tính từ đầu tàu là toa thứ mấy tính từ đuôi tàu lên?
\shortans{4}
\loigiai{
\textbf{Đáp số:} $4$\par
Vị trí từ đuôi tàu lên là: $10 - 7 + 1 = 4$.
\textbf{Đáp số:} $4$.
}
\end{ex}

\begin{ex} %[1D6V24]
Tìm số tự nhiên $m$ có hai chữ số thỏa mãn: chữ số hàng chục hơn chữ số hàng đơn vị là 7 và hàng đơn vị là 1.
\shortans{81}
\loigiai{
\textbf{Đáp số:} $81$\par
Chữ số hàng chục là: $1 + 7 = 8$. Số đó là $81$.
\textbf{Đáp số:} $81$.
}
\end{ex}

\begin{ex} %[1D6V24]
Tính tổng của 4 chục và 5 chục.
\shortans{90}
\loigiai{
\textbf{Đáp số:} $90$\par
Ta có: 4 chục + 5 chục = 9 chục, tức là $90$.
\textbf{Đáp số:} $90$.
}
\end{ex}

\begin{ex} %[1D6V24]
Tính hiệu của 9 chục và 6 chục.
\shortans{30}
\loigiai{
\textbf{Đáp số:} $30$\par
Ta có: 9 chục - 6 chục = 3 chục, tức là $30$.
\textbf{Đáp số:} $30$.
}
\end{ex}

\begin{ex} %[1D6V24]
Số nào cộng với 50 thì bằng 90?
\shortans{40}
\loigiai{
\textbf{Đáp số:} $40$\par
Số đó là: $90 - 50 = 40$.
\textbf{Đáp số:} $40$.
}
\end{ex}

\begin{ex} %[1D6V24]
Số nào trừ đi 20 thì bằng 60?
\shortans{80}
\loigiai{
\textbf{Đáp số:} $80$\par
Số đó là: $60 + 20 = 80$.
\textbf{Đáp số:} $80$.
}
\end{ex}
